\documentclass[10pt]{article}
\usepackage[a4paper,margin=0.72in]{geometry}
\usepackage{amsmath,amssymb,amsthm,mathtools}
\usepackage{graphicx}
\usepackage[hidelinks]{hyperref}
\setlength{\parindent}{0pt}
\setlength{\parskip}{0.45em}
\newtheorem{theorem}{Theorem}
\newtheorem{lemma}[theorem]{Lemma}
\newtheorem{corollary}[theorem]{Corollary}
\newcommand{\R}{\mathbb R}
\newcommand{\norm}[1]{\lVert #1\rVert}
\title{Bounded nilpotency step forces non-degeneracy\\
in direct systems of Carnot groups}
\author{Candidate full solution to Problem P3 of arXiv:2101.03979}
\date{12 August 2026}

\begin{document}
\maketitle

\textbf{Status.} Candidate full solution, likely valid; human review needed.
The answer to Problem P3 is affirmative.

\section*{Source problem}

Moisala and Pasqualetto ask in Problem P3 \cite[p.~3]{MP}:

\begin{center}
\includegraphics[width=0.98\textwidth]{figures/open_problem_crop.png}
\end{center}

In words: if a direct system of infinite-dimensional Carnot groups has a
uniform bound on the nilpotency steps, must it be non-degenerate in the
category of complete metric scalable groups (CMSGs)?

\section*{Proof intuition}

The obstruction in the source example is that conjugation can send elements
tending to the identity a definite distance away.  In a group of fixed
nilpotency class this cannot happen.  A homogeneous component of weighted
degree $r$ can be recovered from an element by a bounded collection of
dilated group words.  Applying that collection to a commutator, and choosing
the two dilation parameters asymmetrically, gives the scale-invariant estimate
\[
 \norm{\exp[U,V]}
 \lesssim_s
 \norm{\exp U}^{p/(p+q)}\norm{\exp V}^{q/(p+q)}
 \qquad (U\in V_p,\ V\in V_q).
\]
Every nontrivial BCH term in $g^{-1}hg$ contains at least one homogeneous
component of $h$ and has total degree at most $s$.  It therefore carries at
least the power $\norm h^{1/s}$.  Crucially, all constants depend only on $s$,
not on the dimension, structure constants, or homogeneous metric.  This
uniformity lets the estimate pass first through the finite-dimensional
filtration of each $G_i$ and then through the infimum pseudometric of the
direct system.

\section*{Homogeneous norms and a uniform collection lemma}

If $G$ is a Carnot group with dilation $\delta$, write
\[
                       \norm g=d(e,g).
\]
Thus $\norm{uv}\leq\norm u+\norm v$, $\norm{u^{-1}}=\norm u$, and
$\norm{\delta_tu}=|t|\norm u$.  No geodesic or sub-Finsler assumption is made.

We record the only nilpotent-collection input.  It is important that its
constant is uniform over the group and the norm.

\begin{lemma}[bounded homogeneous collection]\label{lem:collection}
For every $s\geq1$ there is $A_s<\infty$ with the following property.  Let
$\mathfrak g=V_1\oplus\cdots\oplus V_s$ be any graded nilpotent Lie algebra,
let $G$ be its simply connected group, and equip $G$ with any homogeneous
left-invariant distance.  Every $z\in G$ has an ordered factorization
\begin{equation}\label{eq:ordered}
 z=z_1z_2\cdots z_s,\qquad z_r\in\exp(V_r),
\end{equation}
such that
\begin{equation}\label{eq:factorbound}
                         \norm{z_r}\leq A_s\norm z
                         \qquad(1\leq r\leq s).
\end{equation}
The same assertion holds for a homogeneous pseudonorm by approximation.
\end{lemma}

\begin{proof}
This is the metric form of finite nilpotent collection.  We prove the slightly
stronger assertion that every $z_r$ is a limit of words in dilated copies of
$z$ and their inverses whose \emph{dilation length} (the sum of the absolute
dilation parameters, with integer powers expanded) is bounded only in terms
of $s$.  Such a word $W$ automatically satisfies
\begin{equation}\label{eq:wordbound}
                         \norm{W(z)}
                         \leq \operatorname{len}_\delta(W)\norm z.
\end{equation}

Proceed by induction on $s$.  Class one is immediate.  For the induction
step, let $Z=\exp(V_s)$ and let $\pi:G\to G/Z$.  By induction, every
homogeneous factor $\bar z_r$ of $\pi(z)$, $r<s$, is the limit of words
$W_{r,n}(\pi z)$ of dilation length at most $A_{s-1}$.  Lift the same words to
$G$.  Equation~\eqref{eq:wordbound} puts $W_{r,n}(z)$ in the fixed compact
ball of radius $A_{s-1}\norm z$ (homogeneous balls in a finite-dimensional
graded group are compact).  After taking a subsequence, it converges to an
element
\[
                         w_r=\exp(W_r+C_r),
                         \qquad W_r\in V_r,quad C_r\in V_s,
\]
whose projection is $\bar z_r$.  The top layer is central, so its error is
peeled off by the exact two-dilate identity
\begin{equation}\label{eq:centralpeel}
 \delta_2\!\left(\exp(W_r+W_s)\right)
 \left(\exp(W_r+W_s)\right)^{-2^r}
 =\exp\big((2^s-2^r)W_s\big),
 \quad r<s,
\end{equation}
with $W_s=C_r$.  After the dilation
$\delta_{|2^s-2^r|^{-1/s}}$, this gives $\exp C_r$ at norm at most
\[
 D_{r,s}\norm{w_r},\qquad
 D_{r,s}=\frac{2+2^r}{|2^s-2^r|^{1/s}}.
\]
Consequently $z_r:=w_r\exp(-C_r)\in\exp(V_r)$ and
$\norm{z_r}\leq(1+D_{r,s})A_{s-1}\norm z$.  Substitution of the lifted word
sequences into \eqref{eq:centralpeel} also preserves the stronger bounded-word
assertion, with a new constant depending only on $s$.

Having obtained $z_1,\ldots,z_{s-1}$, set
\[
                         z_s=(z_1\cdots z_{s-1})^{-1}z.
\]
Its image under $\pi$ is the identity, so $z_s\in Z=\exp(V_s)$, and
subadditivity bounds it by
$\norm z+\sum_{r<s}\norm{z_r}$.  Taking the maximum of these finitely many
recursive constants defines $A_s$.  This proves both \eqref{eq:ordered} and
\eqref{eq:factorbound}, with no dimension, structure constant, or metric
equivalence constant.  A homogeneous pseudonorm follows by adding an
arbitrarily small genuine homogeneous norm and letting its coefficient tend
to zero.
\end{proof}

\section*{The uniform conjugation estimate}

\begin{lemma}[weighted commutator estimate]\label{lem:bracket}
If $U\in V_p$, $V\in V_q$, and $p+q\leq s$, then
\begin{equation}\label{eq:bracket}
 \norm{\exp[U,V]}
 \leq B_s
 \norm{\exp U}^{p/(p+q)}
 \norm{\exp V}^{q/(p+q)},
\end{equation}
where $B_s$ depends only on $s$.
\end{lemma}

\begin{proof}
Put $u=\exp U$, $v=\exp V$, $P=\norm u$, and $Q=\norm v$.
For $a,b>0$, the first nonzero homogeneous factor of
$[\delta_a u,\delta_bv]$ is
$\exp(a^pb^q[U,V])$.  Lemma~\ref{lem:collection} and the commutator word give
\[
 (a^pb^q)^{1/(p+q)}\norm{\exp[U,V]}
 \leq A_s\norm{[\delta_a u,\delta_bv]}
 \leq2A_s(aP+bQ).
\]
When $P,Q>0$, choose $a=P^{-1}$ and $b=Q^{-1}$.  This gives
\eqref{eq:bracket} with $B_s=4A_s$.  The cases $P=0$ or $Q=0$ follow by
approximation (and are trivial for a genuine norm).
\end{proof}

Iteration of Lemma~\ref{lem:bracket} gives the corresponding estimate for
every homogeneous Lie monomial.  If a monomial $M$ has total weighted degree
$d\leq s$, and its entries $E_j$ have weights $r_j$, then
\begin{equation}\label{eq:monomial}
 \norm{\exp M(E_1,\ldots,E_m)}
 \leq C_s\prod_j\norm{\exp E_j}^{r_j/d}.
\end{equation}
Rational coefficients in BCH are absorbed into $C_s$, since a scalar on
$V_d$ changes the homogeneous norm by its $1/d$ power.  Finite sums of such
monomials are collected by the same triangular word process as in
Lemma~\ref{lem:collection}; the correction commutators have higher weight and
obey \eqref{eq:monomial} again.  Thus their exponential is bounded by a
constant times the sum of the displayed monomial bounds.

\begin{theorem}[uniform conjugation modulus]\label{thm:conj}
For every $s\geq1$ there is $C_s<\infty$ such that, in every Carnot group of
step at most $s$ with every homogeneous left-invariant distance,
\begin{equation}\label{eq:conj}
 \norm{g^{-1}hg}
 \leq C_s(1+\norm g)^s\norm h^{1/s}
 \qquad\text{whenever }\norm h\leq1.
\end{equation}
The same estimate holds in every infinite-dimensional Carnot group that is
algebraically nilpotent of class at most $s$.
\end{theorem}

\begin{proof}
For a finite-dimensional group, use Lemma~\ref{lem:collection} to write $g$
and $h$ as products of at most $s$ pure homogeneous factors, each bounded by
$A_s\norm g$ or $A_s\norm h$, respectively.  Equivalently, apply the finite
BCH formula to
\[
       \log(g^{-1}hg)=e^{-\operatorname{ad}(\log g)}\log h
\]
and then collect the result into homogeneous factors.  Every nonconstant Lie
monomial that occurs contains at least one homogeneous component coming from
$h$; all remaining entries come from $g$.  If $d\leq s$ is its total weight
and $q\geq1$ is the total weight contributed by its $h$-entries, then
\eqref{eq:monomial} bounds that factor by
\[
        C_s(1+\norm g)^s\norm h^{q/d}
        \leq C_s(1+\norm g)^s\norm h^{1/s},
\]
because $q/d\geq1/s$ and $\norm h\leq1$.  There are only finitely many
universal BCH terms of weight at most $s$.  Subadditivity proves
\eqref{eq:conj}.

Now let $G$ be an infinite-dimensional Carnot group of class at most $s$ and
let $(N_m)$ be a filtration by finite-dimensional Carnot subgroups.  Every
$N_m$ has class at most $s$, so the estimate just proved has the same constant
on all $N_m$.  Approximate $g,h$ by elements of the increasing union
$\bigcup_mN_m$ and pass to the limit.  Multiplication and inversion are
continuous in the already-given CMSG $G$, so \eqref{eq:conj} follows on $G$.
\end{proof}

For comparison, the classical fixed-group estimate has the same exponent
$1/s$, but its stated constant depends on the particular group and metric;
see Magnani \cite[Lemma~3.6]{Mag}.  The uniformity in
Theorem~\ref{thm:conj} is what is needed for a direct system.

\section*{Affirmative solution of P3}

Let
\[
 \mathcal G=\big(\{G_i\}_{i\in I},\{\varphi_{ij}\}_{i\leq j}\big)
\]
be a direct system as in P3, and suppose every $G_i$ has nilpotency class at
most $s$.  Let $(G',\varphi'_i)$ be its direct limit as scalable groups and
let $d'$ be the source's infimum pseudometric.

\begin{theorem}[answer to Problem P3]\label{thm:p3}
The direct system $\mathcal G$ is non-degenerate in the category of CMSGs.
More precisely, for $x=\varphi'_i(x_i)$ and
$M=d_i(e_i,x_i)$, right translation satisfies the uniform-at-$x$ modulus
\begin{equation}\label{eq:rightmod}
 d'(x,yx)\leq C_s(1+M)^s d'(e',y)^{1/s}
 \qquad\text{when }d'(e',y)\leq1,
\end{equation}
with the right side interpreted by an arbitrarily close upper approximation
if the infimum defining $d'$ is not attained.
\end{theorem}

\begin{proof}
Fix $x=\varphi'_i(x_i)$.  Let $y\in G'$ and choose any
$\eta>d'(e',y)$.  By the definition of the infimum pseudometric, there are a
stage $j$ and $a_j,b_j\in G_j$ such that
\[
 \varphi'_j(a_j)=e',\qquad \varphi'_j(b_j)=y,
 \qquad d_j(a_j,b_j)<\eta.
\]
Set $h_j=a_j^{-1}b_j$.  Then $\varphi'_j(h_j)=y$ and, by left invariance,
$\norm{h_j}_j<\eta$.

Choose a common upper stage $k\geq i,j$ and put
\[
 x_k=\varphi_{ik}(x_i),\qquad h_k=\varphi_{jk}(h_j).
\]
The bonding maps are $1$-Lipschitz, so
$\norm{x_k}_k\leq M$ and $\norm{h_k}_k<\eta$.  If $\eta\leq1$,
Theorem~\ref{thm:conj} gives
\begin{align*}
 d'(x,yx)
 &\leq d_k(x_k,h_kx_k)
   =\norm{x_k^{-1}h_kx_k}_k\\
 &\leq C_s(1+M)^s\eta^{1/s}.
\end{align*}
Letting $\eta\downarrow d'(e',y)$ proves \eqref{eq:rightmod}.  In particular,
every right translation on $(G',d')$ is continuous at the identity; this is
condition (C2) of \cite{MP}.  Notice that neither injectivity of the bonding
maps nor attainment of the infimum was used.

It remains to check the two CMSG conditions rather than only (C2).  By source
Lemma 3.10, (C2) makes the metric quotient $G'/Z(d')$ a metric scalable group.
It has class at most $s$: any $(s+1)$-fold commutator has finitely many entries,
hence all are represented at one common stage, where that commutator vanishes.
Source Lemma 3.17 on this quotient therefore upgrades (C2) to
Cauchy-continuity of inversion, which is condition (C3) for the system.
Moreover the first layer generates the quotient as a topological scalable
group by the argument in source equation (4.3); source Lemma 4.7 then gives
the dilation condition (C1).  Thus $\mathcal G$ is non-degenerate as a direct
system of CMSGs.
\end{proof}

\begin{corollary}
Every direct system in P3 has a direct limit in the category of CMSGs.  If the
indexing set is countable, that limit is an infinite-dimensional Carnot group.
\end{corollary}

\begin{proof}
Apply source Theorem 3.15 for existence.  In the countable case apply source
Theorem 4.8 (equivalently Theorem 1.2) to the non-degenerate system.
\end{proof}

\section*{Verification, limitations, and novelty}

Eight focused attempts were recorded.  Counterexample attempts based on
growing structure constants and changing homogeneous metrics fail precisely
because the optimized commutator estimate forces the newly created layer to
shrink.  Classes two and three were collected explicitly; in class three, if
$r=\exp(Y_2+Z_3)$ then
$\delta_2(r)r^{-4}=\exp(4Z_3)$, confirming the uniform layer extraction and
the $1/3$ exponent.  The filtration passage, noninjective bonding maps,
near-minimizers for $d'$, and the source's separate conditions (C1)--(C3) were
then audited.

The theorem does not address P1 or P2.  It resolves exactly P3.  The principal
review point is Lemma~\ref{lem:collection}: its finite collection constants
must depend only on the truncation weight $s$.  The free graded construction
and central peeling identity make this dependence explicit and exclude
dimension, structure constants, and metric equivalence constants.

A bounded novelty search through 12 August 2026 covered the run indexes,
arXiv-id and exact-title searches, the quoted wording of P3, and searches for
bounded-step non-degeneracy and homogeneous conjugation estimates.  It found
the source and its published version, plus the fixed-group estimate in
\cite{Mag}, but no later paper claiming an answer to P3.  Novelty confidence is
moderate: the ingredients are standard nilpotent collection tools, while the
uniform direct-system consequence appears not to have been recorded.

\textbf{Human-review recommendation.} Check the bounded-word proof of
Lemma~\ref{lem:collection} first, then verify the near-minimizing-pair step in
Theorem~\ref{thm:p3}.  Once those two points are accepted, the remaining
implications are explicit source lemmas.

\begin{thebibliography}{9}
\bibitem{MP}
T. Moisala and E. Pasqualetto,
\emph{Direct limits of infinite-dimensional Carnot groups},
Math. Scand. 128 (2022), no.~2, 160--200;
arXiv:2101.03979, doi:10.7146/math.scand.a-132062.

\bibitem{Mag}
V. Magnani,
\emph{Towards differential calculus in stratified groups},
J. Aust. Math. Soc. 95 (2013), no.~1, 76--128;
arXiv:math/0701322.
\end{thebibliography}

\end{document}
